\section{Module Introduction}

Software Engineering has the module code COMP2281 and is led by \hyperref[modulelead-craigstewart]{Dr Craig Stewart}.

\subsection{Assessment}

This module is assessed through a group project, the deadlines for this can be seen in \elemref{tab:Assessments for Software Engineering.}{Table}. Additional formatives are also shown in italics.

\inserttable{c|c|c|c|c|c}{Assessments for Software Engineering.}
{\textbf{Title} & \textbf{Opening} & \textbf{Due} & \textbf{Feedback} & \textbf{Professor} & \textbf{Weight}}
{
    Peer Evaluation 1 & 16th Oct & 19th Nov & 10th Dec & \hyperref[modulelead-craigstewart]{Dr Craig Stewart} & $1.67\%$ \\
    Requirement Documentation & 16th Oct & 19th Nov & 17th Dec & \hyperref[modulelead-craigstewart]{Dr Craig Stewart} & $15\%$ \\
    Peer Evaluation 2 & 16th Oct & 11th Feb & 11th Mar & \hyperref[modulelead-craigstewart]{Dr Craig Stewart} & $1.67\%$ \\
    Product Presentation & Week 16 & - & 18th Mar & \hyperref[modulelead-craigstewart]{Dr Craig Stewart} & $15\%$ \\
    Technical Report & 5th Feb & 18th Mar & 6th May & \hyperref[modulelead-craigstewart]{Dr Craig Stewart} & $30\%$ \\
    Peer Evaluation 3 & 5th Feb & 18th Mar & 6th May & \hyperref[modulelead-craigstewart]{Dr Craig Stewart} & $1.67\%$ \\
    Git Commit & 5th Feb & 29th Apr & 1st Jul & \hyperref[modulelead-craigstewart]{Dr Craig Stewart} & $5\%$ \\
    Project Handover & 5th Feb & 29th Apr & 1st Jul & \hyperref[modulelead-craigstewart]{Dr Craig Stewart} & $10\%$ \\
    Reflective Experience Report & 5th Feb & 6th May & 1st Jul & \hyperref[modulelead-craigstewart]{Dr Craig Stewart} & $20\%$}


\subsection{Timetable}

This module is taught through two one-hour lectures and one two-hour computer class each week as seen below in \elemref{tab:Michaelmas Term Timetable for Software Engineering.}{Table}.

\inserttable{c|c|c|c|c}{Michaelmas Term Timetable for Software Engineering.}
{\textbf{Day} & \textbf{Time} & \textbf{Class Type} & \textbf{Location} & \textbf{Lecturer}}
{
    Tuesday & 09:00 - 10:00 & \colouredText{teal}{Lecture} & TLC033 & \hyperref[lecturer-nellybencomo]{Dr Nelly Bencomo} \\
    Wednesday & 11:00 - 12:00 & \colouredText{teal}{Lecture} & TLC033 & \hyperref[lecturer-nellybencomo]{Dr Nelly Bencomo} \\
    Thursday & 11:00 - 13:00 & \colouredText{orange}{Practical} & CC-0007 & \hyperref[lecturer-gagangeetaujla]{Dr Gagangeet Aujla}}

\separate{7}These lectures are taught jointly by \hyperref[modulelead-craigstewart]{Dr Craig Stewart} and \hyperref[lecturer-nellybencomo]{Dr Nelly Bencomo}.

\section{Project Management}

\subsection{Your Team, Formation, and Leadership}

This lecture was taught by \hyperref[lecturer-nellybencomo]{Dr Nelly Bencomo}.

\separate{7}Belbin (https://www.belbin.com/about/belbin-team-roles) defines nine roles someone can take in a team. These are:
\begin{itemize}
    \item The \underline{Innovator} who solves problems as the \say{ideas guy} of the group.
    \item The \underline{Resource Investigator} who contacts external people/companies to bring additional resources to the group.
    \item The \underline{Coordinator} who guides and controls the team. 
        They know their team's abilities well.
    \item The \underline{Shaper} who is a demanding and confrontational leader who pushes their members to excel.
    \item The \underline{Monitor/Evaluator} who analyses proposed solutions and choices.
    \item The \underline{Team Worker} who facilitates team functions and mediates issues.
    \item The \underline{Implementor} who implements the agreed upon plans.
    \item The \underline{Completer/Finisher} who focuses on details and meeting deadlines.
\end{itemize}

\separate{7}After assigning roles and completing skill evaluation tests, the team can then create a CPD plan which identifies existing skills, areas for development, and a plan for doing that.

\separate{7}Maslow, A. (2016) splits competency into four categories: unconscious incompetence, conscious incompetence, unconscious competence, and finally conscious competence.

\subsection{Team Engagement, Negotiation, and Ideation}

When negotiating, you are more likely to be successful if you have the following factors: legitimacy, confidence, courtesy, adaptation to the other party, rapport, incentives/trade-offs, and research into the bigger picture.

This lecture was taught by \hyperref[lecturer-nellybencomo]{Dr Nelly Bencomo}.

\section{Software Engineering}

\section{Legal, Social, Ethical, and Professional Issues}

This recommended reading for this sub-topic is:
\begin{itemize}
    \item \hyperref[cit:SE Code of Ethics]{The Software Engineering Code of Ethics and Professional Practice, ACM (2024)}.
\end{itemize}

\subsection{Introduction to Software Engineering and Software Engineering Ethics} 

This lecture was taught by \hyperref[modulelead-craigstewart]{Dr Craig Stewart}.

\subsubsection{Introduction to Software Engineering}

\definition{ISP}{An \underline{ill-structured problem}, as defined by \hyperref[cit:The Structure of Ill Structured Problems]{(Simon, 1973)}, is: elusive, unbounded, not understandable, without end conditions. They require orientation and context, and have many stakeholders and many possible solutions.}

\separate{7}\definition{Good Problem Solver}{A \underline{good problem solver} needs the following:
\begin{itemize}
    \item \textbf{Method}: Refers to a formal procedure for accomplishing a goal, typically independent of the tools used.
    \item \textbf{Tool}: An instrument or automated system for accomplishing something in a better way.
    \item \textbf{Procedure}: A combination of tools and techniques to produce a product.
    \item \textbf{Paradigm}: The philosophy or approach for building a project.
\end{itemize}}

\separate{7}\definition{Stakeholders}{For this module, a stakeholder will generally fall into one of three roles:
\begin{itemize}
    \item \textbf{Customer/Client}: The company, organisation, or person who pays for the software system. 
    \item \textbf{Developer}: The company, organisation, or person who is building the software system.
    \item \textbf{User}: The person/people who will actually use the system.
\end{itemize}}

\subsubsection{Introduction to Ethics}

\definition{Software Engineering Code of Ethics}{The software engineering code of ethics is defined by \hyperref[cit:The Dark Side of Software Engineering]{Rost and Glass (2010)} and summarised more thoroughly by the \hyperref[cit:SE Code of Ethics]{Association for Computing Machinery (2024)}.}

\section{Group Project}